\documentclass[12pt]{article}
% The preamble.sty file contains all the packages/options the report uses. it also has defined some custom commands (with a ESB prefix) to make using this template simple.
\usepackage{ESB_preamble} 


%%% Document Variables %%%
\newcommand{\ESBassessment}{Research Project Report}
\newcommand{\ESBtitle}{*project title*}
\newcommand{\ESBauthor}{*author*}
\newcommand{\ESBcoworker}{*co-worker*} % comment out this line if no coworker, if multiple separate names with \\
\newcommand{\ESBsupervisors}{*supervisor1* \\ *supervisor2*} % separate names with \\
\newcommand{\ESBdate}{*report date*} % e.g., October 2001
\newcommand{\ESBstudentid}{*123456789 abcd123*} % or upi or both

\addbibresource{Example_Document/example_references.bib} % Uncomment and edit if using different .bib file name.
% \addbibresource{references.bib} % Uncomment and edit if using different .bib file name.


% If you need to add your own packages/options, add them to preamble.sty or below (but before \begin{document}).







%%% --- %%%


\begin{document}
\pagenumbering{roman} % front matter roman numerals 



%%% Title Page %%%
% The title_page.tex file has the commands to create the title page, it shouldn't need to be changed.
\input{ESB_titlepage.tex} 
%%% --- %%%



%%% Abstract %%%
\ESBfrontsection{Abstract} % Add numberless section and TOC entry
Abstract text.
\clearpage
%%% --- %%%



%%% Declaration %%%
% Replace the unsigned version with a new signed pdf.
\includepdf[pagecommand={
    \addcontentsline{toc}{section}{Declaration}
    \thispagestyle{ESBpagestyle}
    }]{Declaration/Unsigned_Declaration_Print.pdf} 
\clearpage
%%% --- %%%



%%% Table of Contents %%%
\tableofcontents
\listoffigures
\listoftables
\clearpage
%%% --- %%%



%%% Glossary of Terminology %%%
\ESBfrontsection{Glossary of Terminology} % Add numberless section and TOC entry

% Add all terminiology e.g., \ESBaddterminology{$g$}{Gravity}
\ESBaddterminology{Term/Symbol}{Definition} 
\ESBaddterminology{$g$}{Gravity}


\ESBprintglossary % Prints glossary
\clearpage
%%% --- %%%



%%% Abbreviations %%%
\ESBfrontsection{Abbreviations} % Add numberless section and TOC entry

% Add all abbreviations e.g., \ESBaddabbreviation{NZ}{New Zealand}
\ESBaddabbreviation{Symbol}{Name}
\ESBaddabbreviation{NZ}{New Zealand}

\ESBprintabbreviations % Prints abbreviations
\clearpage
%%% --- %%%



%%% Body %%%
\pagenumbering{arabic} % front matter arabic numerals




\section{The .tex File}
A LaTeX file has a .tex file extension.
At the top, it needs the \verb|\documentclass{article}| command to say what type of document it is.
The \verb|article| class is fine for everything ESB.

Underneath are lines of package imports and options for the document. 
The actual formatted text of the document is within the \verb|begin{document} \end{document}| environment.

Case for commands is important. 
Whitespace is not important except when used for new paragraphs.
The \textbackslash is the escape character for commands (i.e., functions). 
The \verb|{}| are the inputs to the commands.
The \verb|[]| are used as options for stuff (can also have key=value pairs).
The \verb|%| is for comments ``Ctrl + /''.

The .tex file structure is similar to a program written in Python or C.
Just like coding, you can get errors when compiling a LaTeX document.
Usually, if you are not doing very peculiar or advanced things, errors are easy to debug and are probably due to a missing bracket, etc.
Overleaf can be very helpful with errors.
Orange warnings should be fixed before finishing the entire document.
Blue warnings are not major issues.
Overleaf and other IDEs/text/editors have many similar shortcuts, such as comment ``Ctrl + /'',  context menu ``Ctrl + space'', and auto-complete ``tab''.



\section{Sections}
Create document sections with the section, subsection, and subsubsection commands:

\section{Section Title}
filler text.
\subsection{Subsection Title}
filler text.
\subsubsection{Subsubsection Title}
filler text.

If you want a numberless section, add a star * after the section commands:

\section*{Numberless Section Title}
filler text.
\subsection*{Numberless Subsection Title}
filler text.

Finer control over sections, i.e., style, numbering, or spacing, is done with the ``titlesec'' package.

\clearpage



\section{Text}
Some normal font text.\\
\textbf{Some bold text.} Note: ``Ctrl + B'' works too.\\
\textit{Some italic text.} Note: Ctrl + I works too.\\
The double backslash \textbackslash\textbackslash~is a line break, i.e., basically latex's \textbackslash n, it is usually unnecessary when writing text.

Empty lines mean to start a new paragraph.

Sentences on each line are a preference, not required though.
But useful cause it makes moving lines to different places easier with ``Alt + arrows'' with the cursor on one or more lines.

A handy tip in overleaf (with different key combos in other IDEs) is duplicating lines with the cursor on one or more lines and pressing ``Alt + shift + arrow keys'' (a quicker copy and paste). 
Useful if you like saving a commented temporary version of something you want to edit or for writing solutions.

The 12pt font size is default and is specified as a document class option.
\href{https://www.overleaf.com/learn/latex/Font_sizes%2C_families%2C_and_styles#Reference_guide}{Other common font sizes (click here) should you need them.}


Note that the \verb.\verb||. command used in this document displays whatever is inside the two symbols (here it is .) as monospaced verbatim (i.e., no commands).

Use \verb|\clearpage| (or \verb|\newpage|) for new page.
\clearpage



\section{Math}
Inline math text is created by enclosing the math with \$ characters.
For example, \verb|$x = 4$| renders as $x = 4$.
Group terms with the curly braces \{ and \}.

To ensure the math is nicely spaced in text and on the same line, you can also surround the inside with curly braces \{ and \} directly after the \$.
Also, add non-breaking spaces \textasciitilde symbol to prevent the math from being "isolated" to the start or end of a line.
For example, \$x~=~4\$ renders as~${x = 4}$.
You can have text in math with the \verb|\text{}| command, for example: $x_\text{start}$.

A numbered equation with the align environment:
\begin{align}
    y = mx + c
\end{align}
An unnumbered equation:
\begin{align*}
    y = ax^2 + bx + c
\end{align*}
We can also align equations with the alignment character \& (while specifying new lines), 
\begin{align*}
    x &= \frac{\pm\sqrt{b^2-4ac}}{2a}\\
    x &= 5 + 10\\
    x &= 15
\end{align*}
A mixture of some common math symbols:
\begin{align*}
    0 &< x \leq 1\\
    y_{23} &= \{[(c_3 \times \beta)]\}\\
    c \div 5 + 1/z &= (\frac{c}{5}) + \left(\frac{c}{5}\right) + z^{-1} \\
    \int_a^b \chi^2 dx &= \sin^2{(\theta)} + C\\
    \frac{dz}{dt} &= \frac{\partial f}{\partial x}\frac{dx}{dt} + ...\\
    \sum_{i=1}^n a_i &= \prod_{j\in \alpha}j
\end{align*}
Set notation is available too, $i\in I \cup \cap \subset \subseteq$.
Matrices too etc., anything ``math'' is available. 
For list of math symbols etc see the \autoref{appendix:other_symbols}, or just search online for a specific one. 

There are many other math environments, look them up if these basic ones here don't suit your needs. 

You can probably find an AI to transform an image (like your handwriting) of an equation to LaTeX.
But obviously, getting AI to solve the said equation \textbf{is against academic integrity}.
\clearpage



\section{Lists}
The enumerate environment is a numbered list.
\begin{enumerate}
    \item 
    \item 
    \item 
\end{enumerate}

The itemize environment is a bullet list.
\begin{itemize}
    \item 
    \item 
    \item 
\end{itemize}

Note these can be nested too:
\begin{enumerate}
    \item Some
        \begin{itemize}
            \item filler
            \item text
        \end{itemize}
    \item to 
        \begin{itemize}
            \item fill
        \end{itemize}
    \item space.
\end{enumerate}

Spacing, numbering, and bullet shape, etc., can be adjusted with the enumitem package.



\section{Files}
A latex document can have multiple files.
The \verb|\input{}| command takes a file path (overleaf can auto-complete (press ``tab'')) and basically copy and pastes the contents of the file into where the input command was called.

Useful for large documents.
Also useful to move the stuff before \verb|\begin{document}| to a separate file to make it look cleaner.

Other slightly different file related commands include: \verb|\include{}|, and \verb|\import{}| (requires package).

\clearpage



\section{Figures}
\label{sec:figure}
To add an image to the document, you must first upload an image to Overleaf.
Then auto-complete (with ``tab'') a \verb|\begin{figure}| environment.
Or you could click and drag the image file into the text and Overleaf will prompt you.

Figure width should be adjusted to your liking (see examples, \autoref{fig:f1} and \autoref{fig:f2}).

\begin{figure}
    \centering
    % 20% width of the text line.
    \includegraphics[width=0.2\linewidth]{fingerprint.png}
    \caption{Fingerprint 1}
    \label{fig:f1}
\end{figure}

\begin{figure}
    \centering
    % Selects the biggest width/height for a given image and keeps the aspect ratio.
    \includegraphics[width=2cm, height=2cm, keepaspectratio]{fingerprint.png}
    \caption{Fingerprint 2}
    \label{fig:f2}
\end{figure}

Latex tries to figure out the best place for the image itself, but you can have more control, with the positional parameters:
\begin{itemize}
    \item \textbf{h: here.} 
    \item t: top of the page. 
    \item b: bottom of the page. 
    \item !: try to force this position.
\end{itemize}
For example, a figure placed right here, otherwise at the bottom or top:
\begin{figure}[!hbt]
    \centering
    \includegraphics[width=1.25cm]{fingerprint.png}
    \caption{Fingerprint 3}
    \label{fig:f3}
\end{figure}

All figures should have a descriptive caption and be referenced in the text like so \autoref{fig:f3}.
Latex will automatically determine the numbering, and if figures have unique labels, then latex will also automatically reference the correct figure numbers in the text.
More description on cross-referencing is found in \autoref{sec:cross-referencing}.
\clearpage



\subsection{Subfigures}
You might also find the need for subfigures, i.e., multiple figures inside a single figure.

Use the \verb|\usepackage{subfig}|.

For example a 1x3 subfigure:
\begin{figure}[!htb]
    \centering
    \begin{subfigure}[t]{0.3\linewidth}
        \centering
        \includegraphics[width=3cm]{fingerprint.png}
        \caption{F1}
        \label{fig:3finger_f1}
    \end{subfigure}
    \hfill % horizontal space fill
    \begin{subfigure}[t]{0.3\linewidth}
        \centering
        \includegraphics[width=3cm]{fingerprint.png}
        \caption{F2}
        \label{fig:3finger_f2}
    \end{subfigure}
    \hfill % horizontal space fill
    \begin{subfigure}[t]{0.3\linewidth}
        \centering
        \includegraphics[width=3cm]{fingerprint.png}
        \caption{F3}
        \label{fig:3finger_f3}
    \end{subfigure}
    \caption{3 Fingers}
    \label{fig:3finger}
\end{figure}

For example a 2x2 subfigure:
\begin{figure}[!htb]
    \centering
    \begin{subfigure}[t]{0.45\linewidth}
        \centering
        \includegraphics[width=.6\linewidth]{fingerprint.png}
        \caption{F1}
        \label{fig:4finger_f1}
    \end{subfigure}
    \hfill % horizontal space fill
    \begin{subfigure}[t]{0.45\linewidth}
        \centering
        \includegraphics[width=.6\linewidth]{fingerprint.png}
        \caption{F2}
        \label{fig:4finger_f2}
    \end{subfigure}
    \vspace{\baselineskip} % vertical space    
    \begin{subfigure}[t]{0.45\linewidth}
        \centering
        \includegraphics[width=.6\linewidth]{fingerprint.png}
        \caption{F3}
        \label{fig:4finger_f3}
    \end{subfigure}
    \hfill % horizontal space fill
    \begin{subfigure}[t]{0.45\linewidth}
        \centering
        \includegraphics[width=.6\linewidth]{fingerprint.png}
        \caption{F4}
        \label{fig:4finger_f4}
    \end{subfigure}
    \caption{4 Fingers}
    \label{fig:4finger}
\end{figure}

Google or AI would easily help you create the latex code for any $m$-by-$n$ subfigure combo.
\clearpage



\subsection{Wrapped figures}
For people who fit the "I need as much space as possible", 
use the \verb|\usepackage{wrapfig}| to wrap text around a figure.

See the following example, where you need to specify what side of the page with R and L and how big of a space latex should leave for the image (while also specifying the size of the image in this space).
You can also specify how many lines the vertical space should have in [].

% draw line
\vspace{\baselineskip}
\hrule

% BANKS - Best Friends % - filler text
When I needed a shoulder, you'd call me, call me
They get a little older, you'd call me, call me, yeah
I've been walking solo lately, lately, uh
I see all your photos, you dress the same, but you're not

\begin{wrapfigure}[8]{L}{0.25\textwidth}
    \centering
    \includegraphics[width=.6\linewidth]{fingerprint.png}
    \caption{Left Finger}
    \label{fig:leftfinger}
\end{wrapfigure}

Tell me to mean what I say
You're not my friend anymore
Don't even mean what you say
While you're standing there begging me for more
There were exceptions within
You told me you did it just to fit in, yeah
But you never fit in anyway

That's why you were my friend
That's why we were best friends
You were my best friend
We never fit in
And that's why you were my friend
That's why we were best friends
You were my best friend
We never fit in

Once, I sent you flowers to say sorry
You said you were feeling like I just wasn't listening
Said I was "lucky for forgiveness"
I didn't feel like you had something to forgive
And now I'm standing alone with no one to stand with
You might also like

\begin{wrapfigure}{R}{0.33\textwidth}
    \centering
    \includegraphics[width=.8\linewidth]{fingerprint.png}
    \caption{Right Finger}
    \label{fig:rightfinger}
\end{wrapfigure}

Tell me to mean what I say
You're not my friend anymore
Don't even mean what you say
While you're standing there begging me for more
There were exceptions within
You told me you did it just to fit in, yeah
But you never fit in anyway

That's why you were my friend
That's why we were best friends
You were my best friend
We never fit in
And that's why you were my friend
That's why we were best friends
You were my best friend
We never fit in

And I got nothing to gain, it's on me
And I got no one to blame, it's on me
And you don't even mean what you said, it's on me
And you don't even mean what you said, it's on me
And I got nothing to gain, it's on me
And I got no one to blame, it's on me
And you don't even mean what you said, it's on me
And you don't even mean what you said, it's on me
\clearpage



\section{Tables}
To make a table, you can auto-complete a \verb|\begin{table}| environment.
You need to manually move the caption and label to the top if you want it at the top.

Note that a tabular environment is like the grid, and the table environment is the whole thing, including caption label.

There \textbf{MUST} be an alignment character for each column.
Column alignment characters:
\begin{itemize}
    \item c: centre
    \item l: left
    \item r: right
\end{itemize}
| pipe characters for vertical lines, can have multiple, and \verb|\hrule| for horizontal lines.
Alignment character \& for new columns and new rows with \verb|\\|.
e.g., see code for this table:

\begin{table}[!hbt]
    \centering
    \caption{Small table}
    \label{tab:small_table}
    \begin{tabular}{l||r|r}
        P2 & Yes & No\\
        \hline
        P2 & 113.5 & 22\\
        P3 & 142.0 & 5 \\
        P4 & 53.0 & 159
    \end{tabular}
\end{table}

Latex tries to figure out the best place for the table itself, but you can have more control, with the positional parameters:
\begin{itemize}
    \item h: here. 
    \item t: top of the page. 
    \item b: bottom of the page. 
    \item !: try to force this position.
\end{itemize}

\textcolor{red}{There are easier ways than manually typing out an entire table's data...}
\begin{itemize}
    \item Use an online latex table generator like:
    \href{https://www.tablesgenerator.com/}{https://www.tablesgenerator.com/}
    where you can copy and paste an Excel table into it and format the table after the data is pasted in.
    \item Latex can also read data directly from an uploaded .csv file into a table. Search it up if you want to know how.
    \item Python also has latex converters, and panda's dataframe has a to\_latex method.
\end{itemize}
% https://pandas.pydata.org/docs/reference/api/pandas.DataFrame.to_latex.html
\clearpage



\section{Cross-Referencing}
\label{sec:cross-referencing}

All figures and tables should have a descriptive caption and be referenced in the text like so \autoref{fig:f3}.
Latex will automatically determine the numbering of everything, and if they have unique labels, then it will also automatically reference the correct figure numbers in the text.

You can label and reference figures/tables/equations/sections (and even custom stuff).
Put the \verb|\label{}| command with the following prefixes and a \textbf{unique} name afterwards:
\begin{itemize}
    \item \textbf{fig:} for figures
    \item \textbf{tab:} for tables
    \item \textbf{sec:} for sections, subsections etc.
    \item \textbf{eq:} for equations
    \item \textbf{appendix:} for appendix sections
\end{itemize}
The label command should be right after the figure or table caption or the section and right after the equation.
See the entire document for labelling examples.

With the label (i.e., the stuff inside the curly brackets \{\},
you can use \verb|\autoref{}| to cross reference, and Overleaf should already know the existing labels by name.

For example, we can reference our figures like so \autoref{fig:4finger}, and even subfigures too, \autoref{fig:4finger_f2}.
Tables as well, like this table \autoref{tab:small_table}
See \autoref{sec:cross-referencing} on cross-sectioning or \autoref{appendix:A} for the ``stuff''. 
\begin{align}
    y = mx + c \label{eq:linear}
\end{align}
Equations can also be referenced as \eqref{eq:linear} or \autoref{eq:linear}.

\textbf{All cross-referenced things are clickable links (the word and number), and you will be linked to that thing. Try it out on any of the references above.}
\clearpage



\section{Referencing Literature}
Quite easy to reference literature with Latex.

Firstly, potential references need to be stored in .bib file.
Many websites that help you create references can also create them in the bibtex format we need.
For example \href{https://zbib.org/}{https://zbib.org/} where you can download bibtex (in the copy to clipboard submenu)\footnote{a lot of the other websites that look the same seem very sus}.

However, you should be using a reference manager like \href{https://www.zotero.org/}{Zotero}, which can also easily create a .bib file (select, rightclick, export), and with their browser extension, easily add things to your literature collection.
Other examples include: Mendeley and EndNote.

The cite commands you should use are:
\begin{itemize}
    \item \verb|\textcite{}|, for a inline text citation. e.g., \textcite{bertsimas_exact_2024}.
    \item \verb|\parencite{}|, for a text citation with brackets. e.g., \parencite{bertsimas_exact_2024}.
    \item \verb|\citeauthor{}|, for only the papers authors. e.g., \citeauthor{bertsimas_exact_2024}.
\end{itemize}
These citation commands, also support multiple citations by separating with a comma, e.g., \parencite{bertsimas_exact_2024, moore_unofficial_2022, liu_using_2022}.

The \verb|\printbibliography| on the next page automatically prints only the references used in the document.

Depends, but, if the bibliography has extra things, it is probably easier for you to edit the fields in your reference manager than adjusting the latex code to not print it.

\clearpage




%%% Acknowledgments %%%
\ESBfrontsection{Acknowledgments} % Add numberless section and TOC entry
Thanks ESB!
\clearpage
%%% --- %%%



%%% References %%%
\printbibliography[
    heading=bibintoc, % Put reference section into table of contents.
    title={\centering References}, % Title of the bibliography.
    ]
\clearpage
%%% --- %%%



%%% Appendicies %%%
\ESBappendix % Start the appendices


\section{Examples of Other Stuff}
\label{appendix:A}
Other interesting and might be useful.
They may require adding the necessary packages to the top of the document - double-check if its already in the preamble.

\subsection{Other Symbols}
\label{appendix:other_symbols}
\begin{itemize}
    \item \href{https://www.cmor-faculty.rice.edu/~heinken/latex/symbols.pdf}{Click to download a 4-page PDF of common latex symbols.}

    \item \href{https://www3.nd.edu/~nmark/UsefulFacts/LaTeX_symbols.pdf}{Click to download a 22-page PDF of the different latex symbols available.}
    \item \href{https://sites.math.washington.edu/~reu/docs/latex_symbols.pdf}{Click to download a 111-page PDF of the different latex symbols available.}
    \item \href{https://tug.ctan.org/info/symbols/comprehensive/symbols-a4.pdf}{Click to download a 481-page PDF of the different latex symbols available.}
\end{itemize}


\subsection{SI Units}
Use the siunitx package for formatted numbers, quantities, units, etc.
For example:
\begin{itemize}
    \item Angles, \ang{50}
    \item Numbers, \num{2000000}, \num{2e10}
    \item Quantities \qty{9.81}{\metre\per\second}, \qty{44}{\joule\per\mole\per\kelvin}
\end{itemize}
\href{http://mirrors.ctan.org/macros/latex/contrib/siunitx/siunitx.pdf}{Lists of numbers/quantities are also supported, for more info, look at the documentation, click here!}


\subsection{Showing Code}
The listings package is a great way to display code (or text) in a nicely formatted style outside of latex's style influence.
It can also be referenced and labelled similarly to a figure.

For example with the \verb|\begin{lstlisting}| environment:

\begin{lstlisting}
import numpy

arr = numpy.array([1, 2, 3, 4, 5])

print(arr)
\end{lstlisting}

We can also tell it what language it is so it does some code highlighting.

\begin{lstlisting}[language=Python]
import numpy

arr = numpy.array([1, 2, 3, 4, 5])

print(arr)
\end{lstlisting}

You can also import from a file with the \verb|\lstinputlisting{}| command.

To give it a caption, provide a caption to the caption option like so:

\begin{lstlisting}[language=Python, caption=Python example, frame=single, numbers=left,]
import numpy

arr = numpy.array([1, 2, 3, 4, 5])

print(arr)
\end{lstlisting}

It also supports colours, line numbering, etc. 
For more customisation, search for some examples of latex lstlisting for the language you want.

\subsection{Algorithms}
The algorithm2e package creates nicely formatted algorithms.
You might use an algorithm to describe an important function you wrote or as pseudocode for methodology.

For example:
\begin{algorithm}
\caption{Latex or Word}
\label{alg:LoW}
\KwIn{Knowledge}
\For{$x \in n^\text{assignments}$}{
    $L \gets 100$\;
    $W \gets \texttt{rand}(0,100)$\;
    \uIf{$L \geq W$}{\nllabel{line:comparison}
        \Return{Use Latex} \nllabel{line:latex}
    }
    \Else{
        continue\;
    }
}
\end{algorithm}
A cross-referencing example: \autoref{alg:LoW} returns \autoref{line:latex} if \autoref{line:comparison} is true.


\subsection{Drawings}
Latex can draw things with the tikz package, for example, flow chart diagrams, graph networks, etc.

tikz draws things very mathematically/programmatically, i.e., by defining what you draw with coordinates. 

Creating things from scratch with tikz is not recommended, try and ``adapt'' something that already exists, i.e., google/AI etc.

There are applications (and online ones too) that can convert their own ``drawings'' into tikz.
For example, \href{https://www.mathcha.io/}{https://www.mathcha.io/} lets you draw stuff with a graphical interface and lets you generate the tikz version, which you can copy paste.

\textbf{For experts only!}\\
\textbf{Note: Tikz can be very slow and easily cause compile issues.}

\subsubsection{Line Plots}
One can use latex to plot graphs directly with the pgfplots package.
However, this is \textbf{not recommended}. 
Just make the graph in python (or equivalent) and save it as a pdf that you upload and show as a figure (\autoref{sec:figure}).

\subsubsection{Electrical diagrams}
The circuitikz package (which uses tikz) creates electrical stuff.

Use \href{https://tikzmaker.com/}{https://tikzmaker.com/} to create electrical circuits with an online GUI and then copy paste the tikz code.




\subsection{To do lists}
The todonotes package creates ``bubbles'' on the side of the text, like visible comments.
Useful for those who like to do lists as it also creates a collated todolist of all the comments throughout the document.

(Overleaf has a inbuilt comment feature too.)
% \usepackage{todonotes}
% \todo{complete this section!}


\subsection{Animation}
PDFs can have animations with the animate package, which is technically just a series of images.
It can only be seem with JavaScript enabled viewers like Adobe's suite.

\clearpage



\section{Other Helpful Resources}
\label{appendix:B}

The overleaf documentation is extremely helpful, it is a lot more detailed than this document.

Other templates.

You can always internet search for things.

AI can be helpful, especially to point you in the right direction, but it can also give you code that doesn't work.
Be careful and remember academic integrity.

%%% --- %%%



\end{document}